% sections/basin.tex

\section{The self-modeling basin}
\label{sec:basin}

This section establishes notation and recalls the results from
Papers~5 and~7 that serve as inputs to the gravitational identification.

\subsection{Self-modeling and quantum mechanics}

Paper~5~\cite{Paper5} argues that a finite-dimensional system carrying a
faithful self-model necessarily has state spaces $M_n(\mathbb{C})^{\mathrm{sa}}$.
The key steps are: (i)~the sequential product structure of the model
forces the state space to be a Euclidean Jordan algebra
(Jordan-von~Neumann-Wigner classification~\cite{JordanVonNeumannWigner1934}),
and (ii)~faithfulness of the model selects the complex matrix algebras
over the real, quaternionic, and spin-factor alternatives. Observers are
therefore C*-subsystems.

\subsection{The exceptional Jordan algebra}

Paper~7~\cite{Paper7} argues that $\hthree$ is the unique arena for
self-modeling. The argument uses non-composability: every simple
Euclidean Jordan algebra except $\hthree$ can be embedded as a tensor
factor in a larger Jordan algebra~\cite{BarnumGraydonWilce2020}. For a
self-modeling basin, this embedding would extend the system, contradicting
the requirement that the model be \emph{of itself}. The exceptional
algebra $\hthree$ has no such embedding and is therefore the unique
candidate.

The algebra $\hthree$ consists of $3 \times 3$ Hermitian matrices over
the octonions $\mathbb{O}$:
\begin{equation}
\label{eq:h3o-element}
X = \begin{pmatrix}
\alpha & x_3 & \bar{x}_2 \\
\bar{x}_3 & \beta & x_1 \\
x_2 & \bar{x}_1 & \gamma
\end{pmatrix}, \quad
\alpha, \beta, \gamma \in \mathbb{R},\;
x_1, x_2, x_3 \in \mathbb{O},
\end{equation}
with $\dim(\hthree) = 3 + 3 \times 8 = 27$. The Jordan product is
$\jp{X}{Y} = \tfrac{1}{2}(XY + YX)$, which is well-defined and
commutative despite the non-associativity of $\mathbb{O}$. The algebra is
simple, formally real, and has degree~3 (the generic minimal polynomial
has degree~3).

The automorphism group is $\Aut(\hthree) = \Ffour$, the 52-dimensional
exceptional Lie group. The structure group (the group preserving the
cubic norm up to scale) is $\Esix$, a real form of $E_6$ with
signature $(-26)$~\cite{Freudenthal1954}.

\subsection{Peirce decomposition}
\label{sec:peirce}

An observer occupying a rank-1 idempotent $E = E_{11}$ (the projector
onto the first diagonal entry, with $\alpha = 1$ and all other entries
zero) induces the Peirce decomposition of $\hthree$ into eigenspaces of
the multiplication operator $L_E(X) = \jp{E}{X}$:
\begin{equation}
\label{eq:peirce-decomposition}
\hthree = \Vone(E) \oplus \Vhalf(E) \oplus \Vzero(E),
\end{equation}
with eigenvalues $1$, $\tfrac{1}{2}$, and $0$ respectively, and
dimensions:
\begin{equation}
\label{eq:peirce-dims}
27 = \underbrace{1}_{\Vone} \oplus \underbrace{16}_{\Vhalf}
\oplus \underbrace{10}_{\Vzero}.
\end{equation}

The Peirce sectors have the following structure:
\begin{itemize}
\item $\Vone \cong \mathbb{R}$: the observer's degree of freedom
  ($\alpha$).

\item $\Vhalf \cong \mathbb{O}^2$: the 16-dimensional interface,
  parametrized by $(x_2, x_3) \in \mathbb{O}^2$. Paper~7 proved that
  $\Vhalf$ carries one generation of Standard Model fermions under
  $\Spin{10} \to \SMgauge$.

\item $\Vzero = \htwo{O}$: the 10-dimensional Peirce complement,
  parametrized by $(\beta, x_1, \gamma)$. This is the observer's
  ``external world'' and is the subject of this paper.
\end{itemize}

The Peirce multiplication rules are~\cite{McCrimmon2004}:
\begin{equation}
\label{eq:peirce-rules}
\begin{aligned}
\Vhalf \circ \Vhalf &\subset \Vone \oplus \Vzero, \\
\Vzero \circ \Vzero &\subset \Vzero, \\
\Vhalf \circ \Vzero &\subset \Vhalf.
\end{aligned}
\end{equation}
The rule $\Vhalf \circ \Vhalf \subset \Vone \oplus \Vzero$ is
physically significant: measurements of the matter sector ($\Vhalf$)
produce outcomes correlated with both the observer ($\Vone$) and the
external world ($\Vzero$). The rule $\Vzero \circ \Vzero \subset
\Vzero$ means $\Vzero$ is a Jordan subalgebra.

\subsection{Observer independence}

Different rank-1 idempotents give different Peirce decompositions, but
they are all $\Ffour$-conjugate. Freudenthal~\cite{Freudenthal1954}
proved that $\Ffour$ acts transitively on the set of rank-1 idempotents
in $\hthree$, with stabilizer $\Spin{9}$. The orbit is the octonionic
projective plane $\mathbb{OP}^2 = \Ffour/\Spin{9}$, a 16-dimensional
compact manifold. Every observer therefore sees the same algebraic
structure, up to an $\Ffour$ automorphism: the Peirce dimensions, the
multiplication rules, and the cubic norm are all observer-independent.

We have verified this computationally for the idempotents $E_{11}$ and
$E_{22}$: the permutation $P = (1,0,2)$ is a Jordan automorphism
mapping $E_{11} \mapsto E_{22}$, and the induced Peirce decompositions
are isomorphic (both have $\Vzero$ of dimension~10 with
$\det_2$ signature $(+,-,-,-)$).
